%% ma: L10-B13-0083
%% lop: 10
%% chuong: 5
%% bai: 13
%% dang: 10.13.10
%% loai: TN4
%% muc: TH
%% dap_an: "B"
%% nguon: "chuyen-de-cac-so-dac-trung-cua-mau-so-lieu-khong-ghep-nhom-toan-10-chuong-trinh-moi.pdf (GV. Phan Nhật Linh) (P03), Bài 02 Dạng 2 Phần I Câu 8"
%% kien_thuc: "Lớp 10 Bài 13 (trung vị, tứ phân vị)"
%% trang_thai: chinh-thuc
%% ghi_chu: "Tự duyệt theo yêu cầu giáo viên 22/10/2026."
\begin{ex}
Các tứ phân vị của mẫu số liệu $12$; $3$; $6$; $15$; $27$; $33$; $31$; $18$; $29$; $54$; $1$; $8$ là
\choice{$Q_1=7;Q_2=17{,}5;Q_3=30$}{\True $Q_1=7;Q_2=16{,}5;Q_3=30$}{$Q_1=7;Q_2=16{,}5;Q_3=30{,}5$}{$Q_1=7{,}5;Q_2=16{,}5;Q_3=30$}
\loigiai{Sắp xếp: $1$; $3$; $6$; $8$; $12$; $15$; $18$; $27$; $29$; $31$; $33$; $54$.

$Q_2=\dfrac{15+18}2=16{,}5$; $Q_1=\dfrac{6+8}2=7$; $Q_3=\dfrac{29+31}2=30$. Chọn B.}
\end{ex}
